% ============================================================================
% LERNBLATT - Physik Klasse 6: Weg-Zeit-Diagramme
% ============================================================================
% Kurze Einführung (1 Seite) vor MT-003
% ============================================================================

\documentclass[11pt,a4paper]{article}

\usepackage[utf8]{inputenc}
\usepackage[T1]{fontenc}
\usepackage[ngerman]{babel}
\usepackage{geometry}
\usepackage{helvet}
\usepackage{tikz}
\usepackage{tcolorbox}
\usepackage{enumitem}
\usepackage{fancyhdr}
\usepackage{xcolor}

\renewcommand{\familydefault}{\sfdefault}

\geometry{left=2cm, right=2cm, top=1.5cm, bottom=1.5cm}

% Farben
\definecolor{physik}{HTML}{2c5aa0}
\definecolor{lightblue}{HTML}{e8f0fa}
\definecolor{gridcolor}{HTML}{c0d0e8}

% Kopfzeile
\pagestyle{fancy}
\fancyhf{}
\fancyhead[L]{\small\textcolor{physik}{\textbf{Physik Klasse 6}} | Weg-Zeit-Diagramme}
\fancyhead[R]{\small Lernblatt}
\renewcommand{\headrulewidth}{0.4pt}

\tcbuselibrary{skins}

% Farbe für Hefter-Hinweis
\definecolor{hinweis}{HTML}{b35c00}
\definecolor{hinweisbg}{HTML}{fff3e0}

% Merksatz-Box
\newtcolorbox{merksatz}{
    colback=lightblue,
    colframe=physik,
    fonttitle=\bfseries,
    title=Merksatz,
    boxrule=1pt,
    arc=2pt
}

% Hefter-Hinweis-Box (AB-Verweis)
\newtcolorbox{hefterhinweis}{
    colback=hinweisbg,
    colframe=hinweis,
    fonttitle=\bfseries\small,
    title={\textcolor{white}{Arbeitsblatt}},
    boxrule=1.5pt,
    arc=2pt,
    left=6pt, right=6pt, top=4pt, bottom=4pt,
    coltitle=white,
    colbacktitle=hinweis
}

\begin{document}

% ============================================================================
% TITEL
% ============================================================================
\begin{center}
    {\Large\bfseries\textcolor{physik}{Weg-Zeit-Diagramme verstehen}}\\[0.3em]
    {\small So stellst du Bewegungen grafisch dar}
\end{center}
\vspace{0.5em}
\hrule
\vspace{1em}

% ============================================================================
% WAS IST EIN s(t)-DIAGRAMM?
% ============================================================================
\textbf{\textcolor{physik}{Was ist ein Weg-Zeit-Diagramm?}}

\vspace{0.5em}
Ein \textbf{Weg-Zeit-Diagramm} (kurz: s(t)-Diagramm) zeigt, wie sich ein Körper bewegt.

\begin{itemize}[leftmargin=*, itemsep=2pt]
    \item Die \textbf{waagerechte Achse} (x-Achse) zeigt die \textbf{Zeit t} in Sekunden (s).
    \item Die \textbf{senkrechte Achse} (y-Achse) zeigt den \textbf{Weg s} in Metern (m).
\end{itemize}

\vspace{0.5em}
\begin{hefterhinweis}
\textbf{Kasten 1:} Ergänze die Achsenbeschriftung (Zeit, Weg, s, m)
\end{hefterhinweis}

\vspace{0.8em}

% ============================================================================
% BEISPIEL
% ============================================================================
\textbf{\textcolor{physik}{Beispiel: Ein Auto fährt gleichförmig}}

\vspace{0.5em}
\begin{minipage}[t]{0.45\textwidth}
\textbf{Messwerte:}

\vspace{0.3em}
\begin{tabular}{|c|c|}
\hline
\textbf{Zeit t (s)} & \textbf{Weg s (m)} \\
\hline
0 & 0 \\
2 & 10 \\
4 & 20 \\
6 & 30 \\
8 & 40 \\
\hline
\end{tabular}

\vspace{0.8em}
Das Auto legt in jeder Sekunde\\
den gleichen Weg zurück:\\
\textbf{5 Meter pro Sekunde.}
\end{minipage}
\hfill
\begin{minipage}[t]{0.50\textwidth}
\textbf{Diagramm:}

\vspace{0.3em}
\begin{tikzpicture}[scale=0.55]
    % Gitter
    \draw[gridcolor, very thin] (0,0) grid (9,5);

    % Achsen
    \draw[->, thick] (0,0) -- (9.5,0) node[right] {\small t (s)};
    \draw[->, thick] (0,0) -- (0,5.5) node[above] {\small s (m)};

    % Achsenbeschriftung
    \foreach \x in {0,2,4,6,8}
        \draw (\x,0) node[below] {\small \x};
    \foreach \y in {0,1,2,3,4}
        \draw (0,\y) node[left] {\small \pgfmathparse{int(\y*10)}\pgfmathresult};

    % Datenpunkte
    \foreach \x/\y in {0/0, 2/1, 4/2, 6/3, 8/4}
        \filldraw[physik] (\x,\y) circle (4pt);

    % Linie
    \draw[physik, thick] (0,0) -- (8,4);

    % Beschriftung
    \node[physik, right] at (8.2,4) {\small Auto};
\end{tikzpicture}
\end{minipage}

\vspace{1.2em}

% ============================================================================
% SO LIEST DU AB
% ============================================================================
\textbf{\textcolor{physik}{So liest du Werte ab:}}

\vspace{0.3em}
\begin{enumerate}[leftmargin=*, itemsep=2pt]
    \item Suche den Zeitpunkt auf der \textbf{waagerechten Achse} (z.\,B. t = 4 s).
    \item Gehe \textbf{senkrecht nach oben} bis zur Linie.
    \item Lies den Weg auf der \textbf{senkrechten Achse} ab (hier: s = 20 m).
\end{enumerate}

\vspace{1em}

% ============================================================================
% MERKSATZ
% ============================================================================
\begin{merksatz}
\begin{itemize}[leftmargin=*, itemsep=4pt]
    \item \textbf{Steile Gerade} = schnelle gleichförmige Bewegung
    \item \textbf{Flache Gerade} = langsame gleichförmige Bewegung
    \item \textbf{Waagerechte Linie} = Stillstand (der Körper bewegt sich nicht)
    \item \textbf{Kurve} = ungleichförmige Bewegung (Geschwindigkeit ändert sich)
\end{itemize}
\end{merksatz}

\vspace{0.5em}
\begin{hefterhinweis}
\textbf{Kasten 2:} Übertrage den Merksatz (steil = schnell, flach = langsam)\\
\textbf{Aufgabe 1:} Lies Werte aus dem Diagramm ab
\end{hefterhinweis}

\vspace{0.6em}

% ============================================================================
% INTERPRETATION: SCHNELL VS. LANGSAM (mit Grafik)
% ============================================================================
\textbf{\textcolor{physik}{Wie erkennst du schnell und langsam?}}

\vspace{0.3em}
\begin{minipage}[t]{0.55\textwidth}
\noindent Die \textbf{Steigung} der Geraden zeigt die Geschwindigkeit:

\begin{itemize}[leftmargin=*, itemsep=2pt, topsep=3pt]
    \item Je \textbf{steiler} die Gerade, desto \textbf{schneller}\\
          {\small (viel Weg in kurzer Zeit)}
    \item Je \textbf{flacher} die Gerade, desto \textbf{langsamer}\\
          {\small (wenig Weg in gleicher Zeit)}
    \item \textbf{Waagerechte} Linie = Stillstand (v = 0)
\end{itemize}

\vspace{0.2em}
\noindent\textit{Die steilere Gerade gehört zum schnelleren Objekt!}
\end{minipage}
\hfill
\begin{minipage}[t]{0.42\textwidth}
\begin{tikzpicture}[scale=0.48]
    \draw[gridcolor, very thin] (0,0) grid (8,5);
    \draw[->, thick] (0,0) -- (8.5,0) node[right] {\small $t$};
    \draw[->, thick] (0,0) -- (0,5.5) node[above] {\small $s$};
    % Schnell (steil, rot)
    \draw[red!70!black, very thick] (0,0) -- (4,4.5);
    \node[red!70!black, right] at (4.2,4.5) {\small Radfahrer};
    \node[red!70!black, below right] at (4.2,4) {\tiny (schnell)};
    % Langsam (flach, blau)
    \draw[physik, very thick] (0,0) -- (8,3);
    \node[physik, right] at (8.2,3) {\small Fußgänger};
    \node[physik, below right] at (8.2,2.5) {\tiny (langsam)};
    % Stillstand
    \draw[gray, very thick] (0,1) -- (8,1);
    \node[gray, right] at (8.2,1) {\small Baum};
    \node[gray, below right] at (8.2,0.5) {\tiny (Stillstand)};
\end{tikzpicture}
\end{minipage}

\vspace{0.8em}

% ============================================================================
% VIER BEISPIELE
% ============================================================================
\textbf{\textcolor{physik}{Vier verschiedene Bewegungen:}}

\vspace{0.5em}
\begin{center}
\begin{tikzpicture}[scale=0.42]
    % Diagramm 1: Langsam (gleichförmig)
    \begin{scope}[shift={(0,0)}]
        \draw[gridcolor, very thin] (0,0) grid (5,4);
        \draw[->, thick] (0,0) -- (5.5,0) node[right] {\tiny t};
        \draw[->, thick] (0,0) -- (0,4.5) node[above] {\tiny s};
        \draw[physik, thick] (0,0) -- (5,2);
        \node[below] at (2.5,-0.8) {\small \textbf{langsam}};
        \node[below] at (2.5,-1.5) {\tiny (gleichförmig)};
    \end{scope}

    % Diagramm 2: Schnell (gleichförmig)
    \begin{scope}[shift={(7,0)}]
        \draw[gridcolor, very thin] (0,0) grid (5,4);
        \draw[->, thick] (0,0) -- (5.5,0) node[right] {\tiny t};
        \draw[->, thick] (0,0) -- (0,4.5) node[above] {\tiny s};
        \draw[physik, thick] (0,0) -- (5,4);
        \node[below] at (2.5,-0.8) {\small \textbf{schnell}};
        \node[below] at (2.5,-1.5) {\tiny (gleichförmig)};
    \end{scope}

    % Diagramm 3: Stillstand
    \begin{scope}[shift={(14,0)}]
        \draw[gridcolor, very thin] (0,0) grid (5,4);
        \draw[->, thick] (0,0) -- (5.5,0) node[right] {\tiny t};
        \draw[->, thick] (0,0) -- (0,4.5) node[above] {\tiny s};
        \draw[physik, thick] (0,2) -- (5,2);
        \node[below] at (2.5,-0.8) {\small \textbf{Stillstand}};
        \node[below] at (2.5,-1.5) {\tiny (keine Bewegung)};
    \end{scope}

    % Diagramm 4: Ungleichförmig (Kurve)
    \begin{scope}[shift={(21,0)}]
        \draw[gridcolor, very thin] (0,0) grid (5,4);
        \draw[->, thick] (0,0) -- (5.5,0) node[right] {\tiny t};
        \draw[->, thick] (0,0) -- (0,4.5) node[above] {\tiny s};
        % Kurve: erst langsam, dann schneller (Beschleunigung)
        \draw[physik, thick] (0,0) .. controls (2,0.5) and (3,2) .. (5,4);
        \node[below] at (2.5,-0.8) {\small \textbf{ungleichförmig}};
        \node[below] at (2.5,-1.5) {\tiny (v ändert sich)};
    \end{scope}
\end{tikzpicture}
\end{center}

\vspace{0.5em}
\begin{hefterhinweis}
\textbf{Kasten 3:} Ergänze die vier Bewegungsarten unter den Diagrammen\\
\textbf{Aufgabe 2:} Zeichne das Diagramm für den Jogger
\end{hefterhinweis}

\end{document}
